% ============================================================================
%  PRECÁLCULO — Sesión 15
%  Semana 15: 16 al 21 de noviembre de 2026
%  3.4 Fracciones algebraicas · 3.5 Funciones trigonométricas
%  Repaso integrador de la Unidad 3 y simulacro del Examen 3
%  Duración prevista: 3 horas
% ============================================================================
\documentclass[aspectratio=169,10pt]{beamer}

\usepackage[utf8]{inputenc}
\IfFileExists{spanish.ldf}{\usepackage[spanish,es-noshorthands]{babel}}{}
\usepackage{amsmath,amssymb}
\usepackage{booktabs}
\usepackage{tikz}
\usepackage{cancel}

\usetheme{Uniminuto}

\uniasignatura{Precálculo}
\uniprofesor{Andrés Felipe Puerta Velez}
\unicorreo{andres.puerta-v@uniminuto.edu.co}
\uninrc{MATE1161}
\unisemana{Sesión 15 — Semana 15 · 16 al 21 de noviembre de 2026}

\newcommand{\listatiny}{\setbeamerfont{itemize/enumerate body}{size=\tiny}\tiny}

\begin{document}

\portada

\begin{frame}{Orden del día}
  \begin{enumerate}
    \item Llamado a lista y revisión de los ejercicios de factorización
    \item \azul{Tema 3.4: fracciones algebraicas}
    \item \azul{Tema 3.5: funciones trigonométricas}
    \item Repaso integrador de la Unidad 3
    \item Simulacro del Examen 3 y recomendaciones
  \end{enumerate}

  \vspace{0.25cm}
  {\footnotesize\color{uniGris} Última sesión de clase del semestre. Lunes 16
  fue festivo (Independencia de Cartagena). Hoy cerramos los dos temas que
  faltan y dejamos la unidad lista: el \destacar{Examen 3 (30\,\%)} se presenta
  la próxima semana en Aulas Virtuales.}
\end{frame}

% ===========================================================================
\seccion{3.4 Fracciones algebraicas}
% ===========================================================================

\begin{frame}{Qué es una fracción algebraica}
  \scriptsize
  \begin{block}{Definición}
    \scriptsize
    Es un cociente de dos polinomios, \ $\dfrac{P(x)}{Q(x)}$, \ con
    $Q(x)\neq 0$. Se comporta \azul{exactamente igual} que una fracción
    numérica: las mismas reglas de simplificar, sumar, multiplicar y dividir.
    Lo único nuevo es que \destacar{el denominador puede anularse para ciertos
    valores de $x$}.
  \end{block}

  \vspace{0.06cm}
  \begin{center}\tiny
  \begin{tabular}{@{}p{3.6cm}p{4.0cm}p{4.2cm}@{}}
    \toprule
    \textbf{Fracción} & \textbf{Se anula el denominador si\dots} &
    \textbf{Dominio} \\
    \midrule
    $\dfrac{x+1}{x-3}$ & $x-3=0 \Rightarrow x=3$ &
      Todos los reales \azul{salvo} $3$ \\
    \addlinespace
    $\dfrac{5}{x^{2}-4}$ & $x^{2}-4=0 \Rightarrow x=\pm 2$ &
      Salvo $2$ y $-2$ \\
    \addlinespace
    $\dfrac{x}{x^{2}+1}$ & Nunca: $x^{2}+1>0$ &
      \destacar{Todos los reales} \\
    \bottomrule
  \end{tabular}
  \end{center}

  \vspace{0.04cm}
  {\tiny\color{uniGris} \destacar{Primer reflejo de todo el tema:} antes de
  operar, factorice el denominador y anote qué valores quedan prohibidos. Esas
  restricciones \azul{sobreviven a la simplificación} y hay que declararlas en
  la respuesta.}
\end{frame}

\begin{frame}{El error que cuesta más puntos en el examen}
  \scriptsize
  \begin{columns}[T,onlytextwidth]
    \begin{column}{0.48\textwidth}
      \begin{alertblock}{Mal}
        \tiny
        \[ \frac{x+3}{x} \;\neq\; \frac{\cancel{x}+3}{\cancel{x}} \;=\; 3 \]
        \[ \frac{x^{2}+9}{x+3} \;\neq\; x+3 \]
        Se tacharon términos que estaban \destacar{sumando}.
      \end{alertblock}
    \end{column}
    \begin{column}{0.48\textwidth}
      \begin{exampleblock}{Bien}
        \tiny
        \[ \frac{3(x+3)}{3} = x+3 \]
        \[ \frac{(x+3)(x-1)}{(x+3)} = x-1 \]
        Se cancelaron \destacar{factores completos}.
      \end{exampleblock}
    \end{column}
  \end{columns}

  \vspace{0.15cm}
  \begin{block}{La regla, en una frase}
    \scriptsize
    \destacar{Solo se cancelan factores, nunca sumandos.} Y para tener factores
    hay que \azul{factorizar primero} —de ahí que la sesión 14 fuera sobre
    factorización—. Si el numerador o el denominador tienen un signo $+$ o $-$
    sin paréntesis alrededor, todavía no se puede cancelar nada.
  \end{block}

  \vspace{0.08cm}
  {\tiny\color{uniGris} Prueba rápida de que está mal: reemplace $x=1$ en
  $\frac{x+3}{x}$: da $4$, no $3$.
  \destacar{Un contraejemplo numérico basta para detectar el error.}}
\end{frame}

\begin{frame}{Simplificar: factorizar, cancelar, declarar}
  \scriptsize
  \begin{exampleblock}{Ejemplo desarrollado}
    \scriptsize
    \[ \frac{x^{2}-9}{x^{2}+7x+12} \]
    \textbf{Paso 1 — factorizar arriba y abajo:}
    \[ \frac{(x-3)(x+3)}{(x+3)(x+4)} \]
    \textbf{Paso 2 — restricciones (del denominador original):}
    $x\neq -3$ \ y \ $x\neq -4$.

    \vspace{0.05cm}
    \textbf{Paso 3 — cancelar el factor común y declarar:}
    \[ \frac{x-3}{x+4}, \qquad x\neq -3,\ x\neq -4 \]
  \end{exampleblock}

  \vspace{0.08cm}
  {\tiny\color{uniGris} Note que $x\neq -3$ \destacar{no se ve} en el resultado
  final, pero sigue valiendo: la función original no existe en $-3$, y
  simplificar no crea puntos nuevos. Por eso el paso 2 va \azul{antes} del
  paso 3.}
\end{frame}

\begin{frame}{Ejercicio guiado 1 — simplificación}
  \begin{alertblock}{En parejas, 10 minutos}
    \scriptsize
    Simplifique y declare las restricciones:
    \begin{enumerate}\setlength{\itemsep}{-0.02cm}
      \item $\dfrac{x^{2}-4}{x^{2}-4x+4}$
      \item $\dfrac{2x^{2}+6x}{x^{2}-9}$
      \item $\dfrac{x^{2}-x-6}{x^{2}-9}$
      \item $\dfrac{3x^{2}-27}{6x+18}$
    \end{enumerate}
  \end{alertblock}

  \vspace{0.04cm}
  {\tiny\color{uniGris} Pista para el 2: saque primero el factor común
  $2x$. Pista para el 4: el $3$ del numerador y el $6$ del denominador también
  se simplifican.}
\end{frame}

\begin{frame}{Ejercicio guiado 1 — solución}
  \scriptsize
  \begin{columns}[T,onlytextwidth]
    \begin{column}{0.49\textwidth}
      \begin{block}{Puntos 1 y 2}
        \tiny
        \textbf{1.} $\dfrac{(x-2)(x+2)}{(x-2)^{2}}
        = \destacar{\dfrac{x+2}{x-2}}$, \ $x\neq 2$

        \vspace{0.12cm}
        \textbf{2.} $\dfrac{2x(x+3)}{(x-3)(x+3)}
        = \destacar{\dfrac{2x}{x-3}}$, \ $x\neq \pm 3$
      \end{block}
    \end{column}
    \begin{column}{0.49\textwidth}
      \begin{block}{Puntos 3 y 4}
        \tiny
        \textbf{3.} $\dfrac{(x-3)(x+2)}{(x-3)(x+3)}
        = \destacar{\dfrac{x+2}{x+3}}$, \ $x\neq \pm 3$

        \vspace{0.12cm}
        \textbf{4.} $\dfrac{3(x-3)(x+3)}{6(x+3)}
        = \destacar{\dfrac{x-3}{2}}$, \ $x\neq -3$
      \end{block}
    \end{column}
  \end{columns}

  \vspace{0.1cm}
  {\tiny\color{uniGris} En el \destacar{4} el resultado ya no tiene
  denominador con $x$, pero la restricción $x\neq -3$ se mantiene: la
  expresión de partida no estaba definida ahí. \azul{Verificación con
  $x=0$:} original $=\frac{-27}{18}=-1{,}5$; simplificada
  $=\frac{-3}{2}=-1{,}5$ \checkmark}
\end{frame}

\begin{frame}{Multiplicar y dividir}
  \scriptsize
  \begin{columns}[T,onlytextwidth]
    \begin{column}{0.48\textwidth}
      \begin{block}{Multiplicación}
        \tiny
        \[ \frac{A}{B}\cdot\frac{C}{D} = \frac{A\cdot C}{B\cdot D} \]
        \destacar{Nunca se desarrollan los productos.} Se factoriza todo, se
        cancela en cruz y solo al final se escribe el resultado.
      \end{block}
    \end{column}
    \begin{column}{0.48\textwidth}
      \begin{block}{División}
        \tiny
        \[ \frac{A}{B}\div\frac{C}{D}
           = \frac{A}{B}\cdot\frac{D}{C} \]
        Se invierte \azul{la segunda} fracción y se multiplica. Igual que con
        números.
      \end{block}
    \end{column}
  \end{columns}

  \vspace{0.12cm}
  \begin{exampleblock}{Ejemplo}
    \tiny
    \[ \frac{x^{2}-1}{x+3}\cdot\frac{x^{2}+6x+9}{x-1}
     = \frac{(x-1)(x+1)}{(x+3)}\cdot\frac{(x+3)^{2}}{(x-1)}
     = (x+1)(x+3) = \destacar{x^{2}+4x+3} \]
    Restricciones: $x\neq -3$, \ $x\neq 1$.
  \end{exampleblock}

  \vspace{0.06cm}
  {\tiny\color{uniGris} Si hubiéramos multiplicado los polinomios de entrada,
  habríamos tenido que factorizar un producto de grado 4 para volver a
  cancelar. \destacar{Factorizar primero ahorra la mitad del trabajo.}}
\end{frame}

\begin{frame}{Ejercicio guiado 2 — producto y cociente}
  \scriptsize
  \begin{alertblock}{En parejas, 10 minutos}
    \scriptsize
    \begin{enumerate}\setlength{\itemsep}{0.06cm}
      \item $\dfrac{x^{2}-9}{x+2}\cdot\dfrac{x+2}{x-3}$
      \item $\dfrac{x+2}{x^{2}-9}\div\dfrac{x^{2}+4x+4}{x-3}$
    \end{enumerate}
  \end{alertblock}

  \vspace{0.1cm}
  \begin{block}{Solución}
    \tiny
    \textbf{1.} $\dfrac{(x-3)(x+3)}{(x+2)}\cdot\dfrac{(x+2)}{(x-3)}
    = \destacar{x+3}$, \ $x\neq -2,\ x\neq 3$

    \vspace{0.1cm}
    \textbf{2.} $\dfrac{(x+2)}{(x-3)(x+3)}\cdot\dfrac{(x-3)}{(x+2)^{2}}
    = \dfrac{1}{(x+3)(x+2)}
    = \destacar{\dfrac{1}{x^{2}+5x+6}}$, \ $x\neq \pm 3,\ x\neq -2$
  \end{block}

  \vspace{0.06cm}
  {\tiny\color{uniGris} En el \destacar{2}, $x\neq 3$ viene del denominador
  original y $x\neq -2$ de que la fracción divisora no puede ser cero. Ambas
  desaparecen del resultado y ambas siguen valiendo.}
\end{frame}

\begin{frame}{Sumar y restar: el mínimo común denominador}
  \scriptsize
  \begin{block}{El procedimiento, en cuatro pasos}
    \tiny
    \textbf{1.} Factorizar todos los denominadores. \quad
    \textbf{2.} Formar el m.c.d.: cada factor distinto, con el \azul{mayor}
    exponente que aparezca. \quad
    \textbf{3.} Amplificar cada fracción para que llegue a ese denominador.
    \quad
    \textbf{4.} Sumar \destacar{solo} los numeradores y simplificar.
  \end{block}

  \vspace{0.1cm}
  \begin{exampleblock}{Ejemplo desarrollado}
    \tiny
    \[ \frac{x}{x-2}-\frac{4}{x^{2}-4}
     = \frac{x}{x-2}-\frac{4}{(x-2)(x+2)} \]
    m.c.d. $=(x-2)(x+2)$. Amplifico la primera por $(x+2)$:
    \[ = \frac{x(x+2)-4}{(x-2)(x+2)}
       = \frac{x^{2}+2x-4}{x^{2}-4}, \qquad x\neq \pm 2 \]
    El numerador $x^{2}+2x-4$ no factoriza en enteros: la respuesta se deja
    así.
  \end{exampleblock}

  \vspace{0.06cm}
  {\tiny\color{uniGris} \destacar{Error clásico:} olvidar el paréntesis al
  restar. $-\,(x^{2}-3x)$ es $-x^{2}+3x$, no $-x^{2}-3x$. Cuando reste, escriba
  el paréntesis \azul{siempre}, aunque le parezca innecesario.}
\end{frame}

\begin{frame}{Ejercicio guiado 3 — suma y resta}
  \scriptsize
  \begin{alertblock}{Individual, 10 minutos}
    \tiny
    Efectúe: \quad
    \textbf{1.} $\dfrac{1}{x}+\dfrac{1}{x+1}$ \qquad
    \textbf{2.} $\dfrac{3}{x+1}+\dfrac{2}{x-1}$ \qquad
    \textbf{3.} $\dfrac{1}{x^{2}-x}-\dfrac{1}{x^{2}-1}$
  \end{alertblock}

  \vspace{0.04cm}
  \begin{block}{Solución}
    \tiny
    \textbf{1.} m.c.d. $=x(x+1)$: \quad
    $\dfrac{(x+1)+x}{x(x+1)} = \destacar{\dfrac{2x+1}{x^{2}+x}}$,
    \ $x\neq 0,-1$

    \vspace{0.08cm}
    \textbf{2.} m.c.d. $=(x+1)(x-1)$: \quad
    $\dfrac{3(x-1)+2(x+1)}{x^{2}-1}
    = \destacar{\dfrac{5x-1}{x^{2}-1}}$, \ $x\neq \pm 1$

    \vspace{0.05cm}
    \textbf{3.} $\dfrac{1}{x(x-1)}-\dfrac{1}{(x-1)(x+1)}$; m.c.d.
    $=x(x-1)(x+1)$: \quad
    $\dfrac{(x+1)-x}{x(x-1)(x+1)} = \destacar{\dfrac{1}{x^{3}-x}}$,
    \ $x\neq 0,\pm 1$
  \end{block}
\end{frame}

\begin{frame}{Fracciones complejas: la fracción de fracciones}
  \scriptsize
  \begin{columns}[T,onlytextwidth]
    \begin{column}{0.47\textwidth}
      \begin{block}{Método}
        \tiny
        \textbf{1.} Convertir el numerador en \azul{una sola} fracción.
        \textbf{2.} Hacer lo mismo con el denominador.
        \textbf{3.} Dividir: multiplicar por la inversa.
      \end{block}
    \end{column}
    \begin{column}{0.49\textwidth}
      \begin{exampleblock}{Ejemplo}
        \tiny
        \[ \frac{\ \dfrac{1}{x}+\dfrac{1}{y}\ }
                {\ \dfrac{1}{x}-\dfrac{1}{y}\ }
         = \frac{\ \frac{y+x}{xy}\ }{\ \frac{y-x}{xy}\ }
         = \frac{y+x}{xy}\cdot\frac{xy}{y-x}
         = \destacar{\frac{x+y}{y-x}} \]
      \end{exampleblock}
    \end{column}
  \end{columns}

  \vspace{0.12cm}
  \begin{block}{Otro ejemplo, paso a paso}
    \tiny
    \[ \frac{\ \frac{x}{x+1}-1\ }{\ 1+\frac{1}{x}\ }
     = \frac{\ \frac{x-(x+1)}{x+1}\ }{\ \frac{x+1}{x}\ }
     = \frac{-1}{x+1}\cdot\frac{x}{x+1}
     = \destacar{\frac{-x}{(x+1)^{2}}}, \qquad x\neq 0,\ x\neq -1 \]
  \end{block}
\end{frame}

\begin{frame}{Ecuaciones racionales y soluciones extrañas}
  \scriptsize
  \begin{block}{El procedimiento}
    \tiny
    \textbf{1.} Anotar las restricciones. \quad
    \textbf{2.} Multiplicar toda la ecuación por el m.c.d. \quad
    \textbf{3.} Resolver la ecuación polinómica que queda. \quad
    \textbf{4.} \destacar{Descartar} las soluciones que violen alguna
    restricción.
  \end{block}

  \vspace{0.1cm}
  \begin{columns}[T,onlytextwidth]
    \begin{column}{0.47\textwidth}
      \begin{exampleblock}{Con solución}
        \tiny
        \[ \frac{2}{x+1}+\frac{1}{x-1}=\frac{3}{x^{2}-1} \]
        Restricciones: $x\neq\pm1$. Por $(x+1)(x-1)$:
        \[ 2(x-1)+(x+1)=3 \Rightarrow 3x-1=3 \]
        \[ x=\destacar{\tfrac{4}{3}} \quad \text{(válida)} \]
      \end{exampleblock}
    \end{column}
    \begin{column}{0.49\textwidth}
      \begin{alertblock}{Sin solución}
        \tiny
        \[ \frac{1}{x-2}+\frac{1}{x+2}=\frac{4}{x^{2}-4} \]
        Restricciones: $x\neq\pm2$. Por $(x-2)(x+2)$:
        \[ (x+2)+(x-2)=4 \Rightarrow 2x=4 \Rightarrow x=2 \]
        Pero $2$ \destacar{está prohibido}: es una \azul{solución extraña}.
        La ecuación \destacar{no tiene solución}.
      \end{alertblock}
    \end{column}
  \end{columns}

  \vspace{0.05cm}
  {\tiny\color{uniGris} El paso 4 no es opcional. Multiplicar por algo que
  \azul{podría ser cero} puede fabricar soluciones que la ecuación original
  jamás tuvo.}
\end{frame}

\begin{frame}{Ejercicio guiado 4 — ecuaciones racionales}
  \scriptsize
  \begin{alertblock}{Individual, 12 minutos}
    \tiny
    Resuelva: \quad
    \textbf{1.} $\dfrac{3}{x-1}=\dfrac{2}{x+4}$ \qquad
    \textbf{2.} $\dfrac{x}{x-2}=\dfrac{4}{x-2}+3$ \qquad
    \textbf{3.} $\dfrac{1}{x}+\dfrac{1}{x+1}=\dfrac{5}{6}$
  \end{alertblock}

  \vspace{0.04cm}
  \begin{block}{Solución}
    \tiny
    \textbf{1.} $3(x+4)=2(x-1) \Rightarrow 3x+12=2x-2 \Rightarrow
    x=\destacar{-14}$ \ (válida: $\neq 1,\ \neq -4$)

    \vspace{0.08cm}
    \textbf{2.} Por $(x-2)$: $x = 4+3(x-2) \Rightarrow x = 3x-2
    \Rightarrow x=\destacar{1}$ \ (válida). Comprobación:
    $\frac{1}{-1}=-1$ y $\frac{4}{-1}+3=-1$ \checkmark

    \vspace{0.05cm}
    \textbf{3.} Por $6x(x+1)$: $6(x+1)+6x=5x(x+1) \Rightarrow
    5x^{2}-7x-6=0 \Rightarrow (5x+3)(x-2)=0$
    \[ x=\destacar{-\tfrac{3}{5}} \quad\text{o}\quad x=\destacar{2}
       \qquad \text{(las dos válidas)} \]
  \end{block}
\end{frame}

% ===========================================================================
\seccion{3.5 Funciones trigonométricas}
% ===========================================================================

\begin{frame}{El triángulo rectángulo: tres cocientes con nombre}
  \scriptsize
  \begin{columns}[T,onlytextwidth]
    \begin{column}{0.44\textwidth}
      \vspace{0.15cm}
      \begin{center}
      \begin{tikzpicture}[scale=0.95,every node/.style={font=\tiny}]
        \draw[uniAzul,very thick] (0,0)--(3.4,0)--(3.4,2.1)--cycle;
        \draw[uniGris] (3.15,0) rectangle (3.4,0.25);
        \draw[uniTitulo,thick] (0.75,0) arc (0:31.7:0.75);
        \node[uniTitulo] at (1.0,0.22) {$\theta$};
        \node[below,uniGris] at (1.7,-0.05) {adyacente};
        \node[right,uniGris] at (3.45,1.05) {opuesto};
        \node[above left,uniGris] at (1.9,1.05) {hipotenusa};
      \end{tikzpicture}
      \end{center}
    \end{column}
    \begin{column}{0.53\textwidth}
      \begin{block}{Las tres razones}
        \tiny
        \[ \operatorname{sen}\theta = \frac{\text{op}}{\text{hip}}
           \qquad
           \cos\theta = \frac{\text{ady}}{\text{hip}}
           \qquad
           \tan\theta = \frac{\text{op}}{\text{ady}} \]
        \vspace{-0.1cm}
        \begin{center}
          \destacar{SOH \quad CAH \quad TOA}
        \end{center}
        Y sus inversas: $\csc=\frac{1}{\operatorname{sen}}$, \
        $\sec=\frac{1}{\cos}$, \ $\cot=\frac{1}{\tan}$.
      \end{block}
    \end{column}
  \end{columns}

  \vspace{0.1cm}
  \begin{alertblock}{Lo que hace útil a la trigonometría}
    \tiny
    Estas razones \destacar{no dependen del tamaño del triángulo}, solo del
    ángulo. Dos triángulos semejantes con el mismo $\theta$ dan el mismo
    $\operatorname{sen}\theta$. Por eso se puede medir la altura de un edificio
    \azul{sin subirse a él}: basta un ángulo y una distancia.
  \end{alertblock}
\end{frame}

\begin{frame}{Los ángulos notables: la tabla que hay que saber}
  \scriptsize
  \begin{center}
  \begin{tabular}{@{}lccccc@{}}
    \toprule
    & $0^{\circ}$ & $30^{\circ}$ & $45^{\circ}$ & $60^{\circ}$ &
      $90^{\circ}$ \\
    \midrule
    $\operatorname{sen}$ &
      $0$ & $\dfrac{1}{2}$ & $\dfrac{\sqrt{2}}{2}$ &
      $\dfrac{\sqrt{3}}{2}$ & $1$ \\
    \addlinespace
    $\cos$ &
      $1$ & $\dfrac{\sqrt{3}}{2}$ & $\dfrac{\sqrt{2}}{2}$ &
      $\dfrac{1}{2}$ & $0$ \\
    \addlinespace
    $\tan$ &
      $0$ & $\dfrac{\sqrt{3}}{3}$ & $1$ & $\sqrt{3}$ &
      \destacar{no existe} \\
    \bottomrule
  \end{tabular}
  \end{center}

  \vspace{0.12cm}
  \begin{exampleblock}{El truco para no memorizarla}
    \tiny
    Escriba $0,1,2,3,4$; divida cada uno entre $4$; saque la raíz. Obtiene
    $\ 0,\ \frac{1}{2},\ \frac{\sqrt{2}}{2},\ \frac{\sqrt{3}}{2},\ 1$: esa es
    la fila del \azul{seno}. La del \azul{coseno} es \destacar{la misma leída
    al revés}. Y $\tan = \operatorname{sen}/\cos$.
  \end{exampleblock}

  \vspace{0.05cm}
  {\tiny\color{uniGris} $\tan 90^{\circ}$ no existe porque el coseno vale $0$:
  \destacar{no se divide entre cero}, igual que en el tema anterior.}
\end{frame}

\begin{frame}{Ejercicio guiado 5 — resolver triángulos}
  \scriptsize
  \begin{alertblock}{En parejas, 12 minutos \ (use calculadora en modo grados)}
    \scriptsize
    \begin{enumerate}\setlength{\itemsep}{0.04cm}
      \item Los catetos de un triángulo rectángulo miden $5$ y $12$. Halle la
            hipotenusa y las tres razones del ángulo opuesto al cateto $5$.
      \item La hipotenusa mide $20$ y un ángulo agudo es $35^{\circ}$. Halle
            el cateto opuesto.
      \item Desde $40$ m de distancia, la punta de una torre se ve con un
            ángulo de elevación de $32^{\circ}$. ¿Qué altura tiene?
    \end{enumerate}
  \end{alertblock}

  \vspace{0.1cm}
  {\scriptsize\color{uniGris} En el 3, dibuje el triángulo antes de calcular:
  la distancia horizontal es el \azul{cateto adyacente} y la altura, el
  \azul{opuesto}. Eso decide qué razón usar.}
\end{frame}

\begin{frame}{Ejercicio guiado 5 — solución}
  \scriptsize
  \begin{columns}[T,onlytextwidth]
    \begin{column}{0.49\textwidth}
      \begin{block}{Punto 1}
        \tiny
        $h=\sqrt{5^{2}+12^{2}}=\sqrt{169}=\destacar{13}$
        \[ \operatorname{sen}\theta=\tfrac{5}{13}\approx 0{,}385 \qquad
           \cos\theta=\tfrac{12}{13}\approx 0{,}923 \]
        \[ \tan\theta=\tfrac{5}{12}\approx 0{,}417
           \ \Rightarrow\ \theta \approx \destacar{22{,}62^{\circ}} \]
      \end{block}
    \end{column}
    \begin{column}{0.49\textwidth}
      \begin{block}{Puntos 2 y 3}
        \tiny
        \textbf{2.} $\operatorname{sen}35^{\circ}=\frac{op}{20}
        \Rightarrow op = 20\operatorname{sen}35^{\circ}
        \approx \destacar{11{,}47}$

        \vspace{0.12cm}
        \textbf{3.} $\tan 32^{\circ}=\frac{h}{40}
        \Rightarrow h = 40\tan 32^{\circ}
        \approx \destacar{25{,}0\ \text{m}}$
      \end{block}
    \end{column}
  \end{columns}

  \vspace{0.1cm}
  {\tiny\color{uniGris} \destacar{Control de sentido común:} en el 2, el cateto
  ($11{,}47$) tiene que ser menor que la hipotenusa ($20$) \checkmark. En el 3,
  como $32^{\circ}<45^{\circ}$, la altura debe ser \azul{menor} que la
  distancia horizontal: $25<40$ \checkmark. Si le da al revés, revisó mal la
  razón.}
\end{frame}

\begin{frame}{Grados y radianes}
  \scriptsize
  \begin{block}{La equivalencia madre}
    \scriptsize
    \[ 180^{\circ} = \pi \ \text{rad}
       \qquad\Longrightarrow\qquad
       \text{rad} = \text{grados}\cdot\frac{\pi}{180}
       \qquad
       \text{grados} = \text{rad}\cdot\frac{180}{\pi} \]
  \end{block}

  \vspace{0.1cm}
  \begin{center}
  \begin{tabular}{@{}lcccccccc@{}}
    \toprule
    Grados & $30^{\circ}$ & $45^{\circ}$ & $60^{\circ}$ & $90^{\circ}$ &
      $135^{\circ}$ & $180^{\circ}$ & $270^{\circ}$ & $360^{\circ}$ \\
    \midrule
    Radianes & $\frac{\pi}{6}$ & $\frac{\pi}{4}$ & $\frac{\pi}{3}$ &
      $\frac{\pi}{2}$ & $\frac{3\pi}{4}$ & $\pi$ & $\frac{3\pi}{2}$ &
      $2\pi$ \\
    \bottomrule
  \end{tabular}
  \end{center}

  \vspace{0.12cm}
  \begin{alertblock}{Por qué existen los radianes}
    \tiny
    Un radián es el ángulo que abarca un arco \azul{igual al radio}. No es una
    unidad arbitraria como el grado (¿por qué 360?): es una \destacar{razón
    entre longitudes}, y por eso las fórmulas del cálculo solo funcionan bien
    en radianes. En Cálculo I trabajarán casi siempre en radianes.
  \end{alertblock}

  \vspace{0.05cm}
  {\tiny\color{uniGris} \destacar{Error frecuente en el examen:} dejar la
  calculadora en RAD cuando el problema está en grados.
  $\operatorname{sen}30 = 0{,}5$ en grados, pero $-0{,}988$ en radianes.
  \azul{Verifique el modo antes de empezar.}}
\end{frame}

\begin{frame}{La circunferencia unitaria y los signos}
  \scriptsize
  \begin{columns}[T,onlytextwidth]
    \begin{column}{0.42\textwidth}
      \vspace{0.05cm}
      \begin{center}
      \begin{tikzpicture}[scale=0.66,every node/.style={font=\tiny}]
        \draw[uniGris,->] (-1.7,0)--(1.9,0) node[right]{$x$};
        \draw[uniGris,->] (0,-1.7)--(0,1.9) node[above]{$y$};
        \draw[uniAzul,very thick] (0,0) circle (1.35);
        \draw[uniTitulo,thick] (0,0)--(0.95,0.95);
        \fill[uniTitulo] (0.95,0.95) circle (1.6pt);
        \node[above right,uniTitulo] at (0.95,0.95)
          {$(\cos\theta,\operatorname{sen}\theta)$};
        \node at (0.85,0.42) {\textbf{I}};
        \node at (-0.85,0.85) {\textbf{II}};
        \node at (0.85,-0.9) {\textbf{IV}};
        \node at (-0.85,-0.9) {\textbf{III}};
      \end{tikzpicture}
      \end{center}
    \end{column}
    \begin{column}{0.55\textwidth}
      \begin{block}{Signos por cuadrante}
        \tiny
        \begin{center}
        \begin{tabular}{@{}lccc@{}}
          \toprule
          & $\operatorname{sen}$ & $\cos$ & $\tan$ \\
          \midrule
          I \ ($0$–$90^{\circ}$)    & $+$ & $+$ & $+$ \\
          II \ ($90$–$180^{\circ}$)  & $+$ & $-$ & $-$ \\
          III \ ($180$–$270^{\circ}$) & $-$ & $-$ & $+$ \\
          IV \ ($270$–$360^{\circ}$)  & $-$ & $+$ & $-$ \\
          \bottomrule
        \end{tabular}
        \end{center}
        \vspace{0.05cm}
        Regla: \destacar{Todos, Seno, Tangente, Coseno} (I, II, III, IV).
      \end{block}
    \end{column}
  \end{columns}

  \vspace{0.04cm}
  \begin{exampleblock}{La identidad que sale de aquí}
    \tiny
    Como el punto está sobre la circunferencia de radio $1$, cumple
    $x^{2}+y^{2}=1$: \ $\operatorname{sen}^{2}\theta + \cos^{2}\theta = 1$.
    \destacar{Es el teorema de Pitágoras disfrazado.} De ahí: si
    $\cos\theta=\frac{5}{13}$ en el primer cuadrante, entonces
    $\operatorname{sen}\theta=\sqrt{1-\frac{25}{169}}=\frac{12}{13}$ y
    $\tan\theta=\frac{12}{5}$.
  \end{exampleblock}
\end{frame}

\begin{frame}{Ejercicio guiado 6 — ángulos, radianes e identidad}
  \scriptsize
  \begin{alertblock}{Individual, 10 minutos}
    \scriptsize
    \begin{enumerate}\setlength{\itemsep}{0.04cm}
      \item Convierta $135^{\circ}$ a radianes y $\frac{7\pi}{6}$ a grados.
      \item Calcule sin calculadora:
            $\operatorname{sen}60^{\circ}\cdot\cos 30^{\circ}
             + \operatorname{sen}30^{\circ}\cdot\cos 60^{\circ}$.
      \item Si $\operatorname{sen}\theta=\frac{3}{5}$ y $\theta$ está en el
            primer cuadrante, halle $\cos\theta$ y $\tan\theta$.
    \end{enumerate}
  \end{alertblock}

  \vspace{0.08cm}
  \begin{block}{Solución}
    \tiny
    \textbf{1.} $135\cdot\frac{\pi}{180}
    = \destacar{\frac{3\pi}{4}}$ \ rad; \quad
    $\frac{7\pi}{6}\cdot\frac{180}{\pi}
    = \destacar{210^{\circ}}$

    \vspace{0.08cm}
    \textbf{2.} $\frac{\sqrt{3}}{2}\cdot\frac{\sqrt{3}}{2}
    + \frac{1}{2}\cdot\frac{1}{2}
    = \frac{3}{4}+\frac{1}{4} = \destacar{1}$
    \quad {\color{uniGris}(es $\operatorname{sen}(60+30)
    = \operatorname{sen}90^{\circ}$)}

    \vspace{0.08cm}
    \textbf{3.} $\cos\theta = \sqrt{1-\frac{9}{25}}
    = \sqrt{\frac{16}{25}} = \destacar{\frac{4}{5}}$; \quad
    $\tan\theta = \frac{3/5}{4/5} = \destacar{\frac{3}{4}}$
    \quad {\color{uniGris}(el triángulo 3-4-5)}
  \end{block}
\end{frame}

% ===========================================================================
\seccion{Repaso integrador de la Unidad 3}
% ===========================================================================

\begin{frame}{Mapa de la Unidad 3}
  \scriptsize
  \begin{center}
  \begin{tabular}{@{}clp{7.4cm}@{}}
    \toprule
    \textbf{Tema} & \textbf{Qué se hace} & \textbf{Lo que se evalúa} \\
    \midrule
    3.1 & Operaciones algebraicas &
      Reducir términos semejantes, leyes de exponentes, operar polinomios \\
    \addlinespace
    3.2 & Productos notables &
      $(a\pm b)^{2}$, $(a+b)(a-b)$, $(a\pm b)^{3}$, $(x+a)(x+b)$ \\
    \addlinespace
    3.3 & Factorización &
      Factor común, diferencia de cuadrados, trinomios, agrupación, cubos \\
    \addlinespace
    3.4 & Fracciones algebraicas &
      Simplificar, operar, fracciones complejas, ecuaciones racionales \\
    \addlinespace
    3.5 & Funciones trigonométricas &
      Razones en el triángulo, ángulos notables, radianes, identidad
      pitagórica \\
    \bottomrule
  \end{tabular}
  \end{center}

  \vspace{0.1cm}
  {\tiny\color{uniGris} Los cinco temas están \destacar{encadenados}: sin
  productos notables no hay factorización, y sin factorización no hay
  fracciones algebraicas. Si algo del 3.3 quedó flojo, es lo primero que hay
  que repasar esta semana.}
\end{frame}

\begin{frame}{Simulacro del Examen 3 — parte A (álgebra)}
  \begin{alertblock}{Individual, 20 minutos, sin apuntes}
    \scriptsize
    \begin{enumerate}\setlength{\itemsep}{0.03cm}
      \item Desarrolle: \ $(3x-2)^{2}$
      \item Factorice completamente: \ $2x^{3}-18x$
      \item Factorice completamente: \ $3x^{3}-12x^{2}+12x$
      \item Simplifique y declare restricciones: \
            $\dfrac{x^{2}-16}{x^{2}+8x+16}$
      \item Efectúe: \ $\dfrac{2}{x-3}+\dfrac{5}{x+2}$
      \item Resuelva: \ $\dfrac{5}{x+3}=\dfrac{2}{x}$
    \end{enumerate}
  \end{alertblock}
\end{frame}

\begin{frame}{Simulacro — solución de la parte A}
  \scriptsize
  \begin{columns}[T,onlytextwidth]
    \begin{column}{0.49\textwidth}
      \begin{block}{Puntos 1 a 3}
        \tiny
        \textbf{1.} $9x^{2}-2\cdot 3x\cdot 2+4
        = \destacar{9x^{2}-12x+4}$

        \vspace{0.1cm}
        \textbf{2.} $2x(x^{2}-9)
        = \destacar{2x(x-3)(x+3)}$

        \vspace{0.1cm}
        \textbf{3.} $3x(x^{2}-4x+4)
        = \destacar{3x(x-2)^{2}}$
      \end{block}
    \end{column}
    \begin{column}{0.49\textwidth}
      \begin{block}{Puntos 4 a 6}
        \tiny
        \textbf{4.} $\dfrac{(x-4)(x+4)}{(x+4)^{2}}
        = \destacar{\dfrac{x-4}{x+4}}$, \ $x\neq -4$

        \vspace{0.1cm}
        \textbf{5.} $\dfrac{2(x+2)+5(x-3)}{(x-3)(x+2)}
        = \destacar{\dfrac{7x-11}{x^{2}-x-6}}$

        \vspace{0.1cm}
        \textbf{6.} $5x=2(x+3) \Rightarrow 3x=6
        \Rightarrow x=\destacar{2}$ \ (válida)
      \end{block}
    \end{column}
  \end{columns}

  \vspace{0.08cm}
  {\tiny\color{uniGris} En el \destacar{2} y el \destacar{3}, quedarse en
  $2x(x^{2}-9)$ o $3x(x^{2}-4x+4)$ es factorización \azul{incompleta}: vale
  media respuesta. En el \destacar{4}, olvidar la restricción $x\neq -4$ es el
  descuento más común de todo el examen.}
\end{frame}

\begin{frame}{Simulacro del Examen 3 — parte B (trigonometría)}
  \begin{alertblock}{Individual, 20 minutos, calculadora en modo grados}
    \scriptsize
    \begin{enumerate}\setlength{\itemsep}{0.03cm}
      \setcounter{enumi}{6}
      \item Escriba $\operatorname{sen}$, $\cos$ y $\tan$ de $60^{\circ}$
            \azul{sin calculadora}.
      \item Convierta $150^{\circ}$ a radianes y $\frac{5\pi}{4}$ a grados.
      \item En un triángulo rectángulo la hipotenusa mide $10$ y un ángulo
            agudo es $25^{\circ}$. Halle los dos catetos.
      \item Si $\cos\theta=\frac{5}{13}$ y $\theta$ está en el primer
            cuadrante, halle $\operatorname{sen}\theta$ y $\tan\theta$.
      \item Desde $50$ m de la base de un edificio, su punta se ve con un
            ángulo de elevación de $38^{\circ}$. ¿Qué altura tiene?
      \item ¿Por qué $\tan 90^{\circ}$ no existe?
    \end{enumerate}
  \end{alertblock}
\end{frame}

\begin{frame}{Simulacro — solución de la parte B}
  \scriptsize
  \begin{columns}[T,onlytextwidth]
    \begin{column}{0.49\textwidth}
      \begin{block}{Puntos 7 a 9}
        \tiny
        \textbf{7.} $\operatorname{sen}60^{\circ}=\frac{\sqrt{3}}{2}$, \
        $\cos 60^{\circ}=\frac{1}{2}$, \ $\tan 60^{\circ}=\sqrt{3}$

        \vspace{0.1cm}
        \textbf{8.} $150\cdot\frac{\pi}{180}=\destacar{\frac{5\pi}{6}}$ rad;
        \ $\frac{5\pi}{4}=\destacar{225^{\circ}}$

        \vspace{0.1cm}
        \textbf{9.} op $=10\operatorname{sen}25^{\circ}
        \approx \destacar{4{,}23}$; \
        ady $=10\cos 25^{\circ}\approx \destacar{9{,}06}$
      \end{block}
    \end{column}
    \begin{column}{0.49\textwidth}
      \begin{block}{Puntos 10 a 12}
        \tiny
        \textbf{10.} $\operatorname{sen}\theta
        =\sqrt{1-\frac{25}{169}}=\destacar{\frac{12}{13}}$; \
        $\tan\theta=\destacar{\frac{12}{5}}$

        \vspace{0.1cm}
        \textbf{11.} $h=50\tan 38^{\circ}
        \approx \destacar{39{,}06\ \text{m}}$

        \vspace{0.1cm}
        \textbf{12.} Porque $\tan=\frac{\operatorname{sen}}{\cos}$ y
        $\cos 90^{\circ}=0$: \destacar{división entre cero}.
      \end{block}
    \end{column}
  \end{columns}

  \vspace{0.08cm}
  {\tiny\color{uniGris} Comprobación del \destacar{9} con Pitágoras:
  $4{,}23^{2}+9{,}06^{2} = 17{,}9+82{,}1 = 100 = 10^{2}$ \checkmark. Del
  \destacar{10}: $\left(\frac{12}{13}\right)^{2}+\left(\frac{5}{13}\right)^{2}
  = \frac{144+25}{169} = 1$ \checkmark. \azul{Siempre hay una forma de
  verificarse solo.}}
\end{frame}

\begin{frame}{Errores frecuentes de toda la Unidad 3}
  \scriptsize
  \begin{center}
  \begin{tabular}{@{}p{5.6cm}p{6.4cm}@{}}
    \toprule
    \textbf{Error} & \textbf{Cómo evitarlo} \\
    \midrule
    Cancelar sumandos: $\frac{x+3}{x}=3$ &
      Solo se cancelan \azul{factores}. Factorice primero \\
    \addlinespace
    $(a+b)^{2}=a^{2}+b^{2}$ &
      Falta $2ab$. Compruebe con $a=b=1$: $4\neq 2$ \\
    \addlinespace
    Factorización incompleta &
      Pregúntese si cada factor todavía se puede factorizar \\
    \addlinespace
    No declarar restricciones &
      Anótelas \azul{antes} de simplificar, no después \\
    \addlinespace
    No verificar la solución racional &
      Toda solución que anule un denominador se descarta \\
    \addlinespace
    Perder el signo al restar fracciones &
      Escriba el paréntesis: $-\,(x^{2}-3x) = -x^{2}+3x$ \\
    \addlinespace
    Calculadora en RAD con datos en grados &
      Revise el modo antes del primer cálculo \\
    \bottomrule
  \end{tabular}
  \end{center}
\end{frame}

\begin{frame}{Cierre del curso}
  \scriptsize
  \begin{block}{Lo que se lleva de Precálculo}
    \tiny
    \textbf{Unidad 1:} los campos numéricos y cómo operar con ellos sin
    calculadora. \quad
    \textbf{Unidad 2:} ecuaciones, desigualdades y el plano cartesiano. \quad
    \textbf{Unidad 3:} el lenguaje algebraico completo más la puerta de entrada
    a la trigonometría.

    \vspace{0.08cm}
    Todo esto es el \destacar{prerrequisito operativo de Cálculo I}: allí no se
    vuelve a explicar cómo factorizar ni cómo sumar fracciones algebraicas, se
    da por sabido.
  \end{block}

  \vspace{0.1cm}
  \begin{alertblock}{Para preparar el Examen 3 (30\,\%) — semana 16}
    \tiny
    \textbf{1.} Rehaga el simulacro de hoy \azul{cronometrado} y sin mirar las
    soluciones. \quad
    \textbf{2.} Repase los siete casos de factorización de la sesión 14: son la
    base de casi todos los puntos. \quad
    \textbf{3.} Memorice la tabla de ángulos notables (o el truco de las
    raíces). \quad
    \textbf{4.} Lleve calculadora \destacar{y verifique que esté en grados}.
  \end{alertblock}

  \vspace{0.1cm}
  \begin{center}
    {\footnotesize\color{uniAzul}\textbf{Gracias por el semestre.}
    \ \ Consultas y asesorías: hasta el día anterior al examen.}
  \end{center}
\end{frame}

\begin{frame}{Ejercicios propuestos}
  \vspace{-0.1cm}
  \begin{beamercolorbox}[wd=\textwidth,sep=0.10cm]{block body}
    \centering\scriptsize
    \textbf{1--6:} simplifique y declare restricciones. \quad
    \textbf{7--12:} efectúe. \quad
    \textbf{13--15:} resuelva. \quad
    \textbf{16--20:} trigonometría.
  \end{beamercolorbox}
  \listatiny
  \vspace{0.15cm}
  \begin{columns}[T,onlytextwidth]
    \begin{column}{0.47\textwidth}
      \begin{enumerate}\setlength{\itemsep}{0.16cm}
        \item $\tfrac{x^{2}-25}{x^{2}-10x+25}$
        \item $\tfrac{x^{2}+5x+6}{x^{2}-4}$
        \item $\tfrac{4x^{2}-1}{2x^{2}+x}$
        \item $\tfrac{x^{3}+27}{x^{2}-9}$
        \item $\tfrac{x^{2}-7x+12}{x^{2}-16}$
        \item $\tfrac{x^{3}-8}{x^{2}-4}$
        \item $\tfrac{1}{x+2}+\tfrac{1}{x-2}$
        \item $\tfrac{x}{x^{2}-1}-\tfrac{1}{x-1}$
        \item $\tfrac{x^{2}-9}{x+2}\cdot\tfrac{x+2}{x-3}$
      \end{enumerate}
    \end{column}
    \begin{column}{0.50\textwidth}
      \begin{enumerate}\setlength{\itemsep}{0.16cm}
        \setcounter{enumi}{9}
        \item $\left(1-\tfrac{1}{x}\right)\div(x-1)$
        \item $\tfrac{3}{x+1}+\tfrac{2}{x-1}$
        \item $\tfrac{x+2}{x^{2}-9}\div\tfrac{x^{2}+4x+4}{x-3}$
        \item $\tfrac{5}{x+3}=\tfrac{2}{x}$
        \item $\tfrac{1}{x-2}+\tfrac{1}{x+2}=\tfrac{4}{x^{2}-4}$
        \item $\tfrac{x}{x+2}-\tfrac{3}{x-2}=1$
        \item Calcule $\operatorname{sen}45^{\circ}\cdot\cos45^{\circ}$
        \item Convierta $210^{\circ}$ a radianes
        \item Cateto adyacente $15$ y ángulo $28^{\circ}$: halle la hipotenusa
        \item Catetos $7$ y $24$: halle el ángulo opuesto al $7$
        \item Escalera a $65^{\circ}$ que llega a $6$ m: ¿cuánto mide?
      \end{enumerate}
    \end{column}
  \end{columns}
\end{frame}

\begin{frame}{Respuestas a los ejercicios propuestos}
  \listatiny
  \vspace{0.15cm}
  \begin{columns}[T,onlytextwidth]
    \begin{column}{0.47\textwidth}
      \begin{enumerate}\setlength{\itemsep}{0.16cm}
        \item $\frac{x+5}{x-5}$, \ $x\neq 5$
        \item $\frac{x+3}{x-2}$, \ $x\neq \pm 2$
        \item $\frac{2x-1}{x}$, \ $x\neq 0,\ x\neq -\frac{1}{2}$
        \item $\frac{x^{2}-3x+9}{x-3}$, \ $x\neq \pm 3$
        \item $\frac{x-3}{x+4}$, \ $x\neq \pm 4$
        \item $\frac{x^{2}+2x+4}{x+2}$, \ $x\neq \pm 2$
        \item $\frac{2x}{x^{2}-4}$
        \item $\frac{-1}{x^{2}-1}$
        \item $x+3$, \ $x\neq -2,\ x\neq 3$
        \item $\frac{1}{x}$, \ $x\neq 0,\ x\neq 1$
      \end{enumerate}
    \end{column}
    \begin{column}{0.50\textwidth}
      \begin{enumerate}\setlength{\itemsep}{0.16cm}
        \setcounter{enumi}{10}
        \item $\frac{5x-1}{x^{2}-1}$
        \item $\frac{1}{x^{2}+5x+6}$
        \item $x=2$
        \item \destacar{No tiene solución} ($x=2$ es extraña)
        \item $x=-\frac{2}{5}$
        \item $\frac{\sqrt{2}}{2}\cdot\frac{\sqrt{2}}{2}=\frac{1}{2}$
        \item $\frac{7\pi}{6}$ rad
        \item $h=\frac{15}{\cos 28^{\circ}}\approx 16{,}99$
        \item $\arctan\frac{7}{24}\approx 16{,}26^{\circ}$
        \item $L=\frac{6}{\operatorname{sen}65^{\circ}}\approx 6{,}62$ m
      \end{enumerate}
    \end{column}
  \end{columns}
\end{frame}

\begin{frame}{Referencias}
  \footnotesize
  \begin{itemize}
    \item[{[1]}] Ramos del Olmo, Á. M. (2016). \textit{Matemáticas básicas
      para el acceso a la universidad}. Ediciones Pirámide.
      \texttt{https://elibro.net/es/ereader/uniminuto/49134}
    \item[{[2]}] González Pulido, J. W., Chacón Benavides, J. A. \&
      Fonseca Correa, L. Á. (2023). \textit{Matemática básica} (1.\textsuperscript{a} ed.).
      Universidad Mariana.
      \texttt{https://elibro.net/es/ereader/uniminuto/280647}
    \item[{[3]}] Placencia Valero, J. (2008). \textit{Compendio de matemática
      básica elemental} (1.\textsuperscript{a} ed.). Editorial Tébar Flores.
      \texttt{https://elibro.net/es/ereader/uniminuto/269758}
  \end{itemize}

  \vspace{0.35cm}
  \begin{center}
    {\Large\color{uniAzul}\textbf{¿Preguntas?}}
  \end{center}
\end{frame}

\end{document}
